\subsection{Task-Relevant Information Preservation}
\label{sec:representationanalysis}
We now bound $D_\phi$ by requiring the representation to preserve information from which expert bimanual actions can be reconstructed.

\subsubsection{Semantic Guidance and Decision Representations}
Semantic guidance can encourage an encoder to merge appearance differences while retaining object-level structure relevant to action selection.
However, appending a deterministic semantic map while retaining the original red--green--blue (RGB) input does not itself improve the best coverage constant because the RGB input remains recoverable by projection.
The relevant object is therefore the learned decision representation $Z$, not the number of input modalities.
To control $D_\phi$, this representation must preserve information that supports a common expert-action rule under both demonstration and deployment distributions.
This decision-relevant view is consistent with information-state and behaviorally grounded representation analyses~\cite{subramanian2022ais,zhang2021dbc}.

% \subsubsection{任务信息与专家动作重构}
\subsubsection{Task Information and Expert-Action Reconstruction}
We examine when retained visibility and interaction geometry, together with execution context, suffice to approximate the expert's coordinated actions.
The reconstruction error depends on both what the selected information fails to explain and how accurately it is estimated.

Continue with $\mu=\bar\mu_\pi$ and the demonstration distribution $\nu$.
Write $\pi(h)=p(Z)$ for the same measurable policy readout under both distributions.
We define a task descriptor vector $G\in\R^{d_G}$ to describe the target's exposure and its spatial relationships to the occluder, tools, and goal.
Depending on the prescribed operation, its components can include visible area, separation from the occluder, tool--target displacement, and target--goal proximity.
Semantic masks provide geometric supervision for selected components, rather than a complete description of contact or deformation.
Let $\chi$ denote task-specified proprioceptive and execution context, selected in advance rather than identified with the unrestricted history.
For bimanual coordination, $\chi$ includes both arms' configurations and any available context needed to distinguish exposure, target interaction, and restoration without requiring an explicit phase label.
Assume measurable readouts $r_\chi$ and $r_G$ such that $\chi=r_\chi(Z)$ and $\widehat G=r_G(Z)\in\R^{d_G}$.
For $\rho\in\{\mu,\nu\}$, define $\delta_{G,\rho}=\E_\rho\norm{G-\widehat G}$.
These quantities must be defined under both distributions, with an explicit convention for unavailable components.

% \begin{assumption}[描述量条件下的专家选择]
\begin{assumption}[Descriptor-conditioned expert selection]
\label{ass:expertselection}
There exists a common measurable rule $g_{\mathrm E}$ mapping descriptors and context to $\mathcal A$ such that, for $\rho\in\{\mu,\nu\}$,
\begin{equation}
 \E_\rho\norm{\pi_{\mathrm E}(h)-g_{\mathrm E}(G,\chi)}
 \le\epsilon_{\mathrm{sel},\rho}.
 \label{eq:expertselectionapprox}
\end{equation}
For a finite constant $K_{\mathrm E}\ge0$,
\begin{equation}
 \norm{g_{\mathrm E}(G,\chi)-g_{\mathrm E}(\widehat G,\chi)}
 \le K_{\mathrm E}\norm{G-\widehat G}
 \label{eq:expertselectionsensitivity}
\end{equation}
holds almost surely under both distributions, with finite $\epsilon_{\mathrm{sel},\rho}$ and $\delta_{G,\rho}$.
\end{assumption}
Here, $\epsilon_{\mathrm{sel},\rho}$ measures expert commands unexplained by $(G,\chi)$, including omitted contact or deformation, while $K_{\mathrm E}$ measures sensitivity to descriptor error.
Task relevance alone does not imply this condition: equally successful actions may reflect different expert choices.

Define $\widehat a_{\mathrm E}(Z)=g_{\mathrm E}(\widehat G,\chi)$.
The triangle inequality yields
\begin{equation}
 \E_\rho\norm{\pi_{\mathrm E}(h)-\widehat a_{\mathrm E}(Z)}
 \le b_\rho
 :=\epsilon_{\mathrm{sel},\rho}+K_{\mathrm E}\delta_{G,\rho}.
 \label{eq:expertreconstruction}
\end{equation}
Appendix~\ref{app:expertreconstruction} gives the derivation.
Supervising descriptor readouts targets the estimation term, while the choice of descriptors and context determines how well expert actions can be explained.

% \subsubsection{条件模仿误差差异的约束}
\subsubsection{Bounds on Conditional Imitation-Error Discrepancy}
\label{sec:relationtransfer}
Expert-action reconstruction connects information preservation to the discrepancy in Section~\ref{sec:representationcoverage}.
Because both $\pi(h)=p(Z)$ and $\widehat a_{\mathrm E}(Z)$ depend on the same representation, their squared difference is a common function under $\mu$ and $\nu$.
Expert reconstruction error bounds how far the actual conditional imitation errors depart from this function.

Assume a finite action-set diameter $B_{\mathcal A}=\sup_{a,a'\in\mathcal A}\norm{a-a'}$ under a fixed action scale.
% \begin{proposition}[通过专家动作重构约束条件误差迁移]
\begin{proposition}[Conditional-error transfer through expert reconstruction]
\label{prop:relationtransfer}
Under Assumptions~\ref{ass:representationcoverage} and~\ref{ass:expertselection}, with $b_\rho$ defined in~\eqref{eq:expertreconstruction},
\begin{equation}
 D_\phi\le 2B_{\mathcal A}(b_\mu+C_\phi b_\nu).
 \label{eq:relationdiscrepancybound}
\end{equation}
If Assumption~\ref{ass:actionsensitivity} also holds,
\begin{equation}
 F(\pi)\le F(\pi_{\mathrm E})
 +TL\sqrt{C_\phi R_\nu(\pi)+2B_{\mathcal A}(b_\mu+C_\phi b_\nu)}.
 \label{eq:relationtransferfailure}
\end{equation}
\end{proposition}
The proof is in Appendix~\ref{app:relationtransfer}.
The bound separates the roles of the learning signals and assumptions.
Action imitation targets $R_\nu(\pi)$, representation coverage is summarized by $C_\phi$, descriptor supervision targets $\delta_{G,\rho}$ through $b_\rho$, and $\epsilon_{\mathrm{sel},\rho}$ records task information omitted by the selected descriptors and context.
For a fixed expert baseline, horizon, and action sensitivity, the task-failure bound tightens only when these combined terms decrease.
